\documentclass[12pt,letterpaper]{article}

% Draft aligned with REMEF research-article structure. Compile with pdfLaTeX.
\usepackage[spanish,english]{babel}
\usepackage[T1]{fontenc}
\usepackage[utf8]{inputenc}
\usepackage{newtxtext,newtxmath}
\usepackage[top=2.5cm,bottom=2.5cm,left=3cm,right=3cm]{geometry}
\usepackage{amsmath,amssymb,amsthm}
\usepackage{booktabs}
\usepackage{array}
\usepackage{float}
\usepackage{microtype}
\usepackage{setspace}
\usepackage{hyperref}
\usepackage[round,authoryear]{natbib}
\usepackage{xcolor}

\setstretch{1.0}
\hypersetup{colorlinks=true,linkcolor=black,citecolor=black,urlcolor=blue}
\setlength{\parindent}{0.6cm}
\setlength{\parskip}{0pt}

% Anonymous manuscript by default for double-blind review.
\newif\ifblind
\blindtrue

\newcommand{\Pmeasure}{\mathbb{P}}
\newcommand{\Qmeasure}{\mathbb{Q}}
\newcommand{\E}{\mathbb{E}}
\newcommand{\Var}{\operatorname{Var}}
\newcommand{\RMSE}{\operatorname{RMSE}}
\newcommand{\MAP}{\operatorname{MAP}}

\begin{document}

\selectlanguage{spanish}

\begin{center}
{\Large\bfseries Incertidumbre Bayesiana en Opciones Asiáticas bajo Medida Neutral al Riesgo\par}
\vspace{0.35cm}
{\large Bayesian Uncertainty in Asian Options under Risk-Neutral Valuation\par}
\vspace{0.6cm}
\ifblind
{\normalsize Manuscrito anónimo para revisión por pares\par}
\else
{\normalsize [Nombre del autor] --- [Institución] --- [ORCID]\par}
\fi
\end{center}

\vspace{0.5cm}

\textbf{Resumen.} El objetivo es cuantificar cómo la incertidumbre bayesiana de los parámetros de un movimiento browniano geométrico estimados bajo la medida física se transmite al valor neutral al riesgo de una opción asiática aritmética. Se implementa Metropolis--Hastings en la parametrización $(\mu,\log\sigma)$ y se valora bajo $\mathbb Q$ mediante Monte Carlo con variables antitéticas y control geométrico. Un experimento de 80 réplicas por tamaño muestral compara integración posterior completa, plug-in de media posterior, MAP y máxima verosimilitud contra el precio verdadero $C^{\mathbb Q}(\sigma_0)$. Los resultados muestran que la incertidumbre de $\mu$ no entra al precio libre de arbitraje y que la incertidumbre relevante se concentra en $\sigma$. Full Bayes y plug-in de media posterior presentan errores casi idénticos, mientras el MAP no muestra ventaja sistemática. Se recomienda distinguir estrictamente inferencia bajo $\mathbb P$ y valoración bajo $\mathbb Q$. La principal limitación es el supuesto Black--Scholes. La originalidad consiste en aislar y medir esta propagación para un derivado dependiente de trayectoria mediante un diseño reproducible y no circular.

\textbf{Palabras clave:} inferencia bayesiana; opciones asiáticas; Metropolis--Hastings; valoración neutral al riesgo; incertidumbre paramétrica.

\textbf{Clasificación JEL:} C11; C15; C63; G13; G17; G32.

\vspace{0.35cm}
\selectlanguage{english}
\textbf{Abstract.} This paper quantifies how Bayesian uncertainty in geometric-Brownian-motion parameters estimated under the physical measure propagates into the risk-neutral value of an arithmetic Asian option. Metropolis--Hastings is implemented in the $(\mu,\log\sigma)$ parameterization, while pricing under $\mathbb Q$ uses Monte Carlo with antithetic sampling and a geometric-Asian control variate. An outer experiment with 80 replications per sample size compares full posterior integration, posterior-mean plug-in, MAP plug-in, and maximum likelihood against the external benchmark $C^{\mathbb Q}(\sigma_0)$. The results show that uncertainty in the physical drift $\mu$ does not enter the no-arbitrage price and that relevant parameter uncertainty is concentrated in $\sigma$. Full Bayes and the posterior-mean plug-in have nearly identical pricing errors, whereas MAP provides no systematic advantage. We recommend a strict separation between inference under $\mathbb P$ and valuation under $\mathbb Q$. The main limitation is the Black--Scholes specification. Originality lies in isolating this propagation for a path-dependent derivative with a reproducible, non-circular design.

\textbf{Keywords:} Bayesian inference; Asian options; Metropolis--Hastings; risk-neutral valuation; parameter uncertainty.

\textbf{JEL Classification:} C11; C15; C63; G13; G17; G32.

\selectlanguage{spanish}

\section{Introducción}

La incertidumbre paramétrica y la valoración de derivados son problemas relacionados pero conceptualmente distintos. Los retornos históricos permiten inferir parámetros bajo una medida física $\Pmeasure$, mientras que en un mercado completo tipo Black--Scholes el precio libre de arbitraje de un derivado negociable se obtiene bajo una medida neutral al riesgo $\Qmeasure$. Confundir ambos objetos puede convertir incertidumbre predictiva sobre el rendimiento esperado en una aparente incertidumbre de precio que no corresponde al modelo de no arbitraje.

Este trabajo estudia esa separación en una opción asiática de promedio aritmético. El contrato es útil como laboratorio porque su pago depende de toda la trayectoria y su precio carece, en general, de la fórmula elemental disponible para una opción europea. Al mismo tiempo, el modelo Black--Scholes permite identificar con precisión qué parámetro histórico debe propagarse al precio neutral al riesgo: la volatilidad $\sigma$, no el drift físico $\mu$.

El objetivo general es evaluar si integrar la distribución posterior de $\sigma$ produce mejoras materiales frente a estimadores plug-in cuando se valora una opción asiática. La comparación se realiza contra un benchmark externo, $C^{\Qmeasure}(\sigma_0)$, generado con el parámetro verdadero. Este detalle evita evaluar un estimador usando como función de pérdida las mismas muestras posteriores que lo definen.

Las contribuciones son cuatro. Primero, se formaliza la separación $\Pmeasure/\Qmeasure$ dentro de un flujo bayesiano sencillo. Segundo, se implementa Metropolis--Hastings sobre $(\mu,\eta)$, con $\eta=\log\sigma$, eliminando la frontera artificial $\sigma>0$. Tercero, se propaga únicamente la posterior de $\sigma$ hacia un pricer neutral al riesgo con reducción de varianza. Cuarto, se compara full Bayes, media posterior plug-in, MAP plug-in y MLE mediante bias, MAE, RMSE y cobertura en un experimento repetido.

\section{Estado del arte}

El fundamento de la valoración neutral al riesgo se remonta a \citet{blackscholes1973} y \citet{merton1973}. En un mercado completo, el drift físico del activo no determina el precio de un reclamo contingente replicable; bajo $\Qmeasure$ el rendimiento esperado se reemplaza por la tasa libre de riesgo ajustada por dividendos.

Para opciones asiáticas, \citet{kemnavorst1990} desarrollan un método Monte Carlo con reducción de varianza basado en el promedio geométrico, mientras \citet{gemanyor1993} estudian funcionales exponenciales del browniano asociados al promedio aritmético. Estos trabajos muestran que la dependencia de trayectoria modifica el problema computacional sin alterar el principio de valoración neutral al riesgo.

La literatura bayesiana de derivados enfatiza que la incertidumbre paramétrica puede incorporarse mediante densidades predictivas neutrales al riesgo. \citet{romboutsstentoft2014} encuentran que, cuando existe una cantidad sustancial de datos, las diferencias de pricing entre inferencia clásica y bayesiana pueden ser pequeñas, aunque la integración posterior sigue siendo conceptualmente correcta. Este resultado constituye un rival natural para la hipótesis de que full Bayes debe mejorar de manera material el precio en todo régimen muestral.

Metropolis--Hastings es un mecanismo general para muestrear distribuciones conocidas hasta una constante normalizante \citep{metropolis1953,hastings1970}. Su uso en opciones no constituye por sí mismo una novedad: por ejemplo, \citet{capuozzo2021} lo emplean para evaluar integrales de trayectoria. En este trabajo, en cambio, M--H se utiliza exclusivamente para la inferencia estadística bajo $\Pmeasure$; la valoración bajo $\Qmeasure$ se realiza mediante Monte Carlo estándar con control variate. Esta separación es central para la interpretación económica.

\section{Metodología}

\subsection{Modelo histórico bajo la medida física}

Se supone que el activo sigue bajo $\Pmeasure$
\begin{equation}
 dS_t=\mu S_t\,dt+\sigma S_t\,dW_t^{\Pmeasure},
 \label{eq:p_gbm}
\end{equation}
con $\sigma>0$. Por It\^o,
\begin{equation}
 r_i:=\log\frac{S_{t_i}}{S_{t_{i-1}}}
 \sim
 \mathcal N\!\left[
 \left(\mu-\frac{\sigma^2}{2}\right)\Delta t,
 \sigma^2\Delta t
 \right].
 \label{eq:returns}
\end{equation}

Se mantienen priors débiles para facilitar comparación con el proyecto base:
\begin{equation}
 \mu\sim\mathcal N(0,1),
 \qquad
 \sigma\sim\operatorname{InvGamma}(2,0.1).
\end{equation}
La inferencia produce $p(\mu,\sigma\mid\mathcal D)$, donde $\mathcal D=(r_1,\ldots,r_n)$.

\subsection{Metropolis--Hastings en $(\mu,\log\sigma)$}

Se define $\eta=\log\sigma\in\mathbb R$. Si el prior se especifica sobre $\sigma$, la densidad objetivo en $(\mu,\eta)$ requiere el Jacobiano
\begin{equation}
 p(\mu,\eta\mid\mathcal D)
 \propto
 p(\mathcal D\mid\mu,e^{\eta})
 p(\mu)p(e^{\eta})e^{\eta}.
 \label{eq:jacobian}
\end{equation}

El algoritmo propone
\begin{equation}
 \theta^*=\theta^{(k)}+\varepsilon,
 \qquad
 \varepsilon\sim\mathcal N(0,\Sigma_q),
\end{equation}
con $\theta=(\mu,\eta)$. Al ser simétrica la propuesta, la aceptación es
\begin{equation}
 \alpha=\min\left\{1,
 \frac{p(\theta^*\mid\mathcal D)}{p(\theta^{(k)}\mid\mathcal D)}
 \right\}.
\end{equation}
La implementación usa log-densidades para evitar underflow y reporta tasa de aceptación. Para una versión final se recomienda complementar con múltiples cadenas, $\widehat R$, ESS y error Monte Carlo.

\subsection{Cambio a la medida neutral al riesgo}

En el modelo Black--Scholes completo, el proceso relevante para pricing es
\begin{equation}
 dS_t=(r-q)S_t\,dt+\sigma S_t\,dW_t^{\Qmeasure},
 \label{eq:q_gbm}
\end{equation}
con tasa libre de riesgo $r$ y dividend yield $q$. El parámetro $\mu$ de \eqref{eq:p_gbm} desaparece del precio de no arbitraje. Para una call asiática aritmética discretamente monitoreada,
\begin{equation}
 A_T=\frac1m\sum_{j=1}^{m}S_{t_j},
 \qquad
 C^{\Qmeasure}(\sigma)
 =e^{-rT}\E^{\Qmeasure}\left[(A_T-K)^+\right].
 \label{eq:asian_price}
\end{equation}

Por tanto, el objeto bayesiano relevante es el push-forward
\begin{equation}
 p(C\mid\mathcal D)
 =\int
 \delta_{C^{\Qmeasure}(\sigma)}(C)
 p(\sigma\mid\mathcal D)\,d\sigma,
 \label{eq:pushforward}
\end{equation}
no una integral sobre el drift físico dentro de la dinámica de pricing.

\subsection{Monte Carlo y control variate geométrico}

La opción aritmética se estima simulando \eqref{eq:q_gbm}. Se emplean pares antitéticos y el precio cerrado de la opción asiática geométrica discretamente monitoreada como control variate, siguiendo la idea de \citet{kemnavorst1990}. Si $Y$ es el payoff aritmético descontado, $X$ el geométrico descontado y $\E[X]$ es conocido, se usa
\begin{equation}
 Y_{cv}=Y-\beta\{X-\E[X]\},
 \qquad
 \beta=\frac{\operatorname{Cov}(Y,X)}{\Var(X)}.
\end{equation}
El control reduce el ruido numérico y permite atribuir la dispersión restante principalmente a incertidumbre paramétrica y no al Monte Carlo interno.

\subsection{Estimadores de precio}

Se comparan cuatro reglas:
\begin{align}
 \widehat C_{FB}
 &=\E[C^{\Qmeasure}(\sigma)\mid\mathcal D],\\
 \widehat C_{PM}
 &=C^{\Qmeasure}(\E[\sigma\mid\mathcal D]),\\
 \widehat C_{MAP}
 &=C^{\Qmeasure}(\widehat\sigma_{MAP}),\\
 \widehat C_{MLE}
 &=C^{\Qmeasure}(\widehat\sigma_{MLE}).
\end{align}
En general, $\widehat C_{FB}\neq\widehat C_{PM}$ por la no linealidad de $C^{\Qmeasure}(\sigma)$. Una expansión de segundo orden alrededor de $\bar\sigma=\E[\sigma\mid\mathcal D]$ da
\begin{equation}
 \widehat C_{FB}
 \approx
 C^{\Qmeasure}(\bar\sigma)
 +\frac12 C^{\Qmeasure\prime\prime}(\bar\sigma)
 \Var(\sigma\mid\mathcal D),
 \label{eq:curvature}
\end{equation}
lo que muestra que el valor incremental de la integración posterior depende simultáneamente de incertidumbre y curvatura.

\subsection{Diseño de evaluación no circular}

Se fija $\mu_0=0.08$, $\sigma_0=0.25$, $S_0=K=100$, $r=0.03$, $q=0$ y $T=1$. Para cada $n\in\{63,252,1260\}$ se generan $R=80$ bases históricas independientes bajo $\Pmeasure$. Cada base produce una posterior y los cuatro estimadores anteriores. El benchmark es externo a la posterior:
\begin{equation}
 C_0=C^{\Qmeasure}(\sigma_0).
\end{equation}
Para un método $M$ se reportan
\begin{align}
 \operatorname{Bias}_M&=\frac1R\sum_{j=1}^{R}(\widehat C_{M,j}-C_0),\\
 \operatorname{MAE}_M&=\frac1R\sum_{j=1}^{R}|\widehat C_{M,j}-C_0|,\\
 \RMSE_M&=\sqrt{\frac1R\sum_{j=1}^{R}(\widehat C_{M,j}-C_0)^2}.
\end{align}
También se evalúa la cobertura de los intervalos posteriores de precio al 95\%. A diferencia de una pérdida $L_2$ calculada contra las mismas muestras posteriores, estas métricas no garantizan algebraicamente la superioridad de la media.

\subsection{Reproducibilidad, software y uso de inteligencia artificial}

El código de inferencia, pricing, pruebas y experimento repetido se implementó en Python con NumPy, SciPy y pandas. El repositorio contiene semillas deterministas, pruebas unitarias y un flujo de integración continua. Para preservar el doble ciego, el enlace identificable al repositorio deberá sustituirse por un repositorio anónimo durante el arbitraje.

Se utilizó ChatGPT (OpenAI) como herramienta de apoyo para revisión metodológica, refactorización de código y edición del manuscrito. Las ecuaciones, supuestos, ejecución computacional, verificación de resultados, selección de referencias y responsabilidad científica permanecen bajo control del autor. Esta declaración se incluye para cumplir con las directrices editoriales sobre transparencia en el uso de herramientas de inteligencia artificial.

\section{Resultados y discusión}

\subsection{Corrección del caso base}

Con una trayectoria sintética de 252 retornos diarios, la volatilidad verdadera es $\sigma_0=0.25$. La inferencia transformada produce una media posterior cercana a 0.242 y un intervalo creíble aproximado $[0.222,0.264]$. La posterior de $\mu$ permanece mucho más dispersa, pero esta incertidumbre no se introduce en \eqref{eq:q_gbm}.

La Tabla \ref{tab:baseline} muestra el efecto económico de la separación de medidas. El benchmark neutral al riesgo es aproximadamente 6.42. Los estimadores basados en la volatilidad inferida se ubican alrededor de 6.22--6.24, y la distribución posterior de precios es mucho más estrecha que la obtenida al propagar incorrectamente $\mu$ dentro de las trayectorias de pricing.

\begin{table}[H]
\centering
\caption{Caso base: valoración neutral al riesgo de la call asiática}
\label{tab:baseline}
\begin{tabular}{lcc}
\toprule
Método & Precio & Comentario \\
\midrule
Benchmark $C^{\Qmeasure}(\sigma_0)$ & 6.42 & $\sigma_0=0.25$ \\
Full Bayes & 6.24 & integra $p(\sigma\mid\mathcal D)$ \\
Media posterior plug-in & 6.24 & $C^{\Qmeasure}(\E[\sigma\mid D])$ \\
MAP plug-in & 6.22 & moda posterior de $\sigma$ \\
MLE plug-in & 6.23 & volatilidad MLE \\
Intervalo posterior 95\% & $[5.80,6.73]$ & push-forward de $\sigma$ \\
\bottomrule
\end{tabular}

\vspace{0.15cm}
\footnotesize Fuente: elaboración propia mediante simulación reproducible bajo $\Pmeasure$ y $\Qmeasure$.
\end{table}

El cambio de interpretación es sustantivo. La incertidumbre estadística en $\mu$ sigue siendo real para pronóstico de retornos bajo $\Pmeasure$, pero no constituye por sí misma riesgo de precio Black--Scholes. La distribución posterior de precios refleja incertidumbre sobre la volatilidad y error numérico residual controlado.

\subsection{Experimento de muestreo repetido}

La Tabla \ref{tab:rmse} resume el experimento de 80 réplicas por tamaño muestral. La diferencia entre full Bayes y el plug-in de media posterior es prácticamente nula en los tres regímenes. Para $n=63$, ambos obtienen RMSE cercana a 0.469; para $n=252$, alrededor de 0.245; y para $n=1260$, alrededor de 0.124. El MAP es algo peor en muestras cortas y no presenta una ventaja persistente.

\begin{table}[H]
\centering
\caption{Error de pricing contra el benchmark neutral al riesgo}
\label{tab:rmse}
\begin{tabular}{rrrrr}
\toprule
$n$ & Método & Bias & MAE & RMSE \\
\midrule
63 & Full Bayes & -0.001 & 0.354 & 0.469 \\
63 & Media posterior plug-in & -0.001 & 0.354 & 0.469 \\
63 & MAP plug-in & -0.092 & 0.378 & 0.487 \\
63 & MLE plug-in & -0.041 & 0.358 & 0.470 \\
\addlinespace
252 & Full Bayes & 0.014 & 0.199 & 0.245 \\
252 & Media posterior plug-in & 0.014 & 0.199 & 0.245 \\
252 & MAP plug-in & -0.021 & 0.197 & 0.245 \\
252 & MLE plug-in & 0.003 & 0.199 & 0.244 \\
\addlinespace
1260 & Full Bayes & -0.005 & 0.101 & 0.124 \\
1260 & Media posterior plug-in & -0.005 & 0.101 & 0.124 \\
1260 & MAP plug-in & -0.013 & 0.106 & 0.130 \\
1260 & MLE plug-in & -0.007 & 0.101 & 0.124 \\
\bottomrule
\end{tabular}

\vspace{0.15cm}
\footnotesize Fuente: elaboración propia; 80 réplicas independientes por $n$.
\end{table}

La cobertura empírica de los intervalos posteriores de precio al 95\% fue aproximadamente 0.950 para $n=63$, 0.963 para $n=252$ y 0.925 para $n=1260$. Las tasas medias de aceptación fueron 0.524, 0.395 y 0.222, respectivamente. La caída de aceptación con $n$ indica que una propuesta fija no escala óptimamente con la concentración posterior, por lo que la versión final del estudio debe adaptar $\Sigma_q$ o utilizar una etapa de calibración.

\subsection{Interpretación del resultado negativo}

El hecho de que $\widehat C_{FB}$ y $\widehat C_{PM}$ sean casi iguales no constituye un fallo del enfoque bayesiano. Es una consecuencia empírica de que, en este contrato y rango de volatilidad, $C^{\Qmeasure}(\sigma)$ es localmente cercana a lineal y la posterior de $\sigma$ se concentra con rapidez. En términos de \eqref{eq:curvature}, el término de segundo orden es pequeño.

Este resultado es coherente con \citet{romboutsstentoft2014}, quienes encuentran que el impacto de la incertidumbre paramétrica sobre errores de pricing puede ser menor cuando la información disponible es suficiente. La contribución aplicada aquí consiste en mostrar el mismo mecanismo en un derivado dependiente de trayectoria y, simultáneamente, demostrar por qué la incertidumbre del drift histórico no debe mezclarse con la valoración neutral al riesgo.

Desde una perspectiva de gestión de riesgo, la distribución posterior de $C$ sigue siendo útil como cuantificación de incertidumbre del modelo condicionada a la especificación Black--Scholes. Sin embargo, no debe interpretarse como sustituto de una medida de riesgo de pérdidas de cartera, VaR o Expected Shortfall. Esos objetos requieren una distribución de P\&L y un horizonte de riesgo explícito.

\section{Conclusiones, recomendaciones y consideraciones finales}

Este estudio separa inferencia estadística y valoración financiera en un esquema bayesiano para opciones asiáticas. Bajo $\Pmeasure$, los datos informan sobre $(\mu,\sigma)$ y el drift presenta gran incertidumbre. Bajo $\Qmeasure$, el precio Black--Scholes de la opción depende del drift $r-q$ y de $\sigma$; por ello, la posterior relevante para incertidumbre de pricing es $p(\sigma\mid\mathcal D)$.

El experimento repetido no encuentra una superioridad material de full Bayes frente al plug-in basado en la media posterior para la configuración estudiada. En muestras cortas, el MAP tampoco domina y puede introducir sesgo adicional. La evidencia favorece una conclusión más acotada que una afirmación universal de superioridad bayesiana: integrar la posterior es el procedimiento coherente para representar incertidumbre paramétrica, pero su beneficio cuantitativo depende de la curvatura del funcional de precio y de la dispersión posterior.

Se recomiendan cuatro extensiones antes del envío definitivo. Primero, aumentar el número de réplicas y variar moneyness, madurez y frecuencia de monitoreo. Segundo, reportar diagnósticos modernos de MCMC con múltiples cadenas. Tercero, estudiar datos reales o calibración conjunta a opciones para separar incertidumbre histórica y riesgo de especificación. Cuarto, extender el marco a volatilidad estocástica, donde el paso entre $\Pmeasure$ y $\Qmeasure$ incorpora precios de riesgo adicionales y la incertidumbre paramétrica puede adquirir mayor relevancia económica.

La principal limitación es precisamente la fortaleza que permite aislar el mecanismo: el modelo Black--Scholes es completo y obliga a que $\mu$ desaparezca del precio. Por tanto, los resultados no se extrapolan automáticamente a mercados incompletos, modelos con volatilidad estocástica, saltos o primas de riesgo no observadas.

\section*{Financiamiento y conflictos de interés}

\textbf{Financiamiento:} [Completar por el autor antes del envío.]\\
\textbf{Conflictos de interés:} [Completar por el autor antes del envío.]

\section*{Disponibilidad de datos y código}

Los datos del experimento son sintéticos y completamente reproducibles a partir de las semillas especificadas en el código. El repositorio contiene los módulos de inferencia bajo $\Pmeasure$, pricing bajo $\Qmeasure$, pruebas automatizadas y el experimento de muestreo repetido. Para una evaluación doble ciego deberá proporcionarse una copia anonimizada del repositorio.

\begin{thebibliography}{99}

\bibitem[Black y Scholes(1973)]{blackscholes1973}
Black, F. y Scholes, M. (1973). The pricing of options and corporate liabilities. \textit{Journal of Political Economy}, 81(3), 637--654. doi:10.1086/260062.

\bibitem[Capuozzo et al.(2021)]{capuozzo2021}
Capuozzo, P., Panella, E., Schettini Gherardini, T. y Vvedensky, D. D. (2021). Path integral Monte Carlo method for option pricing. \textit{Physica A: Statistical Mechanics and its Applications}, 581, 126231. doi:10.1016/j.physa.2021.126231.

\bibitem[Geman y Yor(1993)]{gemanyor1993}
Geman, H. y Yor, M. (1993). Bessel processes, Asian options, and perpetuities. \textit{Mathematical Finance}, 3(4), 349--375. doi:10.1111/j.1467-9965.1993.tb00092.x.

\bibitem[Glasserman(2004)]{glasserman2004}
Glasserman, P. (2004). \textit{Monte Carlo Methods in Financial Engineering}. Springer.

\bibitem[Hastings(1970)]{hastings1970}
Hastings, W. K. (1970). Monte Carlo sampling methods using Markov chains and their applications. \textit{Biometrika}, 57(1), 97--109.

\bibitem[Kemna y Vorst(1990)]{kemnavorst1990}
Kemna, A. G. Z. y Vorst, A. C. F. (1990). A pricing method for options based on average asset values. \textit{Journal of Banking \& Finance}, 14(1), 113--129. doi:10.1016/0378-4266(90)90039-5.

\bibitem[Metropolis et al.(1953)]{metropolis1953}
Metropolis, N., Rosenbluth, A. W., Rosenbluth, M. N., Teller, A. H. y Teller, E. (1953). Equation of state calculations by fast computing machines. \textit{The Journal of Chemical Physics}, 21(6), 1087--1092.

\bibitem[Merton(1973)]{merton1973}
Merton, R. C. (1973). Theory of rational option pricing. \textit{The Bell Journal of Economics and Management Science}, 4(1), 141--183. doi:10.2307/3003143.

\bibitem[Rombouts y Stentoft(2014)]{romboutsstentoft2014}
Rombouts, J. V. K. y Stentoft, L. (2014). Bayesian option pricing using mixed normal heteroskedasticity models. \textit{Computational Statistics \& Data Analysis}, 76, 588--605. doi:10.1016/j.csda.2013.06.023.

\end{thebibliography}

\end{document}
